\chapter{Partial Fractions and Trigonometric Integrals}

In the previous chapter, integration by parts showed us how to trade a
stubborn product for a friendlier one by shifting derivatives across factors.
Together with substitution, which runs the chain rule in reverse, it gave us
two powerful ways to undo differentiation. And underneath both sits a single
quiet idea that this chapter turns into an explicit strategy: \emph{we almost
never integrate the integrand we are handed}. We reshape it. We push and pull
on it with algebra and trigonometry until it turns into something we already
recognize -- a power, a logarithm, an arctangent -- and then we read off the
answer.

That reshaping is the real subject here. We will meet two families of
integrands that look hopeless at first and become routine once you know the
trick. The first is \emph{rational functions}, ratios of polynomials like
$\dfrac{3x+5}{x^2-x-2}$. The cure is \emph{partial fractions}: a way of
breaking one ugly fraction into a sum of simple ones, each of which integrates
to a logarithm or an arctangent. The second is integrals built from
\emph{powers and products of sines and cosines}, and then a beautiful move
called \emph{trigonometric substitution} that turns a stubborn algebraic
expression like $\sqrt{a^2 - x^2}$ into a clean trig expression by drawing a
right triangle and reading off its sides.

Nothing in this chapter is new calculus: every integral we evaluate reduces to
standard integrals you already know, using algebra and trigonometry to bridge
the gap. One principle runs through all of it: \textbf{every method in this chapter reshapes the integrand into a
known form.} Partial fractions reshapes a ratio into a sum of standard pieces.
Trig identities reshape a product into a sum. Trig substitution reshapes a
square root into a triangle. Completing the square reshapes a quadratic into
a shifted square. Keep asking: \emph{what known form am I steering toward?}

\chapterline

\section{The Unifying Idea: Reshaping Into Known Forms}
\label{ch24:reshape}

Let us name the small set of integrals we are allowed to treat as ``known.''
Everything in this chapter lands on one of these.

\begin{focusbox}
The standard forms we steer toward:
\[
\int x^n\,\dd x = \frac{x^{n+1}}{n+1}\;(n\neq -1),
\qquad
\int \frac{\dd x}{x} = \ln\abs{x},
\]
\[
\int \frac{\dd x}{x^2 + a^2} = \frac{1}{a}\arctan\!\frac{x}{a},
\qquad
\int \frac{\dd x}{\sqrt{a^2 - x^2}} = \arcsin\!\frac{x}{a},
\]
together with $\int e^x\,\dd x = e^x$, $\int \sin x\,\dd x = -\cos x$,
$\int \cos x\,\dd x = \sin x$, and whatever a single $u$-substitution can
reach from these. \emph{That is the entire target vocabulary.} Each technique
below is a translator from a strange integrand into this small language.
\end{focusbox}

Why these particular forms? Because each is a derivative you have already
computed and reversed. $\ln\abs{x}$ is there because its derivative is $1/x$;
the absolute value is there because $\ln$ is defined only for positive inputs, and
$\dfrac{\dd}{\dd x}\ln\abs{x} = 1/x$ on both sides of zero. The arctangent is
there because $\dfrac{\dd}{\dd x}\arctan x = \dfrac{1}{1+x^2}$, and the
$a$-flavored version is just that rule run through a substitution. The
arcsine is there for the same reason with $\dfrac{1}{\sqrt{1-x^2}}$.

\begin{intuitionbox}
Think of integration as a language-translation problem, not a calculation. You
are fluent in maybe a dozen ``words'' (the standard forms). An arbitrary
integrand is a sentence in a foreign language. You do not invent new grammar
to translate it; you rewrite it, clause by clause, using identities you
already know, until every clause is a word you recognize. Partial fractions,
trig identities, and trig substitution are three dictionaries for three
common foreign tongues.
\end{intuitionbox}

\section{Rational Functions and the Promise of Partial Fractions}
\label{ch24:rational-intro}

A \emph{rational function} is a quotient of polynomials, $R(x) = P(x)/Q(x)$.
We would like to integrate any such thing. Two of them are already easy:
\[
\int \frac{\dd x}{x - 3} = \ln\abs{x-3},
\qquad
\int \frac{\dd x}{(x-3)^2} = -\frac{1}{x-3}.
\]
Both are one substitution ($u = x-3$) away from a standard form. The trouble
starts when the denominator factors into \emph{several} different roots, all
tangled together, as in
\[
\int \frac{3x + 5}{x^2 - x - 2}\,\dd x .
\]
The integrand as written is not any standard form. But watch what happens if
we are clever. The denominator factors: $x^2 - x - 2 = (x-2)(x+1)$. Suppose,
for a moment, that we could write
\[
\frac{3x+5}{(x-2)(x+1)} \;=\; \frac{A}{x-2} + \frac{B}{x+1}
\]
for some constants $A$ and $B$. Then the integral would split into two pieces
we \emph{can} do:
\[
\int \frac{A}{x-2}\,\dd x + \int \frac{B}{x+1}\,\dd x
= A\ln\abs{x-2} + B\ln\abs{x+1}.
\]
If every rational function can be decomposed into a sum
of pieces like $\dfrac{A}{x-r}$ (and a couple of cousins we will meet), then
integrating rational functions is finished: each piece is a logarithm or an
arctangent, and we are done.

\begin{focusbox}
\textbf{Partial fraction decomposition} is the claim that any proper rational
function $P(x)/Q(x)$ (where the degree of the numerator is strictly smaller
than the degree of the denominator, just like a proper numerical fraction like
$3/5$ whose top is smaller than its bottom) can be written as a sum of
``simple'' fractions, one group for each factor of the denominator. It converts
an integral you cannot see into a sum of integrals you can.
\end{focusbox}

Two things must be checked before we trust this. Can we always
\emph{find} the constants $A$, $B$? And, deeper, does the decomposition
always \emph{exist}? We find the constants first -- it is satisfying and
quick -- and then come back and prove the decomposition is guaranteed, because
understanding \emph{why it exists} is more illuminating than any single
computation.

\subsection{Finding the constants}

Start from the proposed identity and clear the denominator by multiplying both
sides by $(x-2)(x+1)$:
\[
3x + 5 \;=\; A(x+1) + B(x-2).
\]
This must hold for \emph{every} $x$. That gives us freedom: we can plug in any
$x$ we like and squeeze out the constants. The smart choices are the values
that kill a factor.

Set $x = 2$: the $B$ term dies, leaving $3(2)+5 = A(2+1)$, so $11 = 3A$ and
$A = \tfrac{11}{3}$.

Set $x = -1$: the $A$ term dies, leaving $3(-1)+5 = B(-1-2)$, so $2 = -3B$ and
$B = -\tfrac{2}{3}$.

\begin{intuitionbox}
This trick of plugging in the root is sometimes called the ``cover-up''
method, and it has a clean meaning. To find the coefficient sitting over
$(x-2)$, you cover up that factor in the original denominator and evaluate
what is left at $x = 2$:
\[
A = \left.\frac{3x+5}{x+1}\right|_{x=2} = \frac{11}{3},
\qquad
B = \left.\frac{3x+5}{x-2}\right|_{x=-1} = \frac{2}{-3} = -\frac{2}{3}.
\]
Why does covering up work? Multiplying by $(x-2)$ removes that factor from the
denominator, and then setting $x = 2$ makes every \emph{other} term vanish,
because each of them still carries a factor of $(x-2)$. You are isolating one
coefficient by evaluating at the precise point where its neighbors disappear.
\end{intuitionbox}

So the decomposition is
\[
\frac{3x+5}{(x-2)(x+1)} = \frac{11/3}{x-2} - \frac{2/3}{x+1},
\]
and the integral falls out:
\[
\int \frac{3x+5}{x^2-x-2}\,\dd x
= \frac{11}{3}\ln\abs{x-2} - \frac{2}{3}\ln\abs{x+1} + C.
\]

\begin{labbox}
Check the decomposition without integrating. Put $\dfrac{11/3}{x-2} -
\dfrac{2/3}{x+1}$ over the common denominator $(x-2)(x+1)$. The numerator is
$\tfrac{11}{3}(x+1) - \tfrac{2}{3}(x-2) = \tfrac{11}{3}x + \tfrac{11}{3}
- \tfrac{2}{3}x + \tfrac{4}{3} = 3x + 5$. It matches. The algebra of putting
fractions over a common denominator is exactly what partial fractions
\emph{undoes}; that is why a check is just running the machine forward.
\end{labbox}

\subsection{Why the decomposition must exist}
\label{ch24:why-pf-exists}

It is one thing to find $A$ and $B$ when they are promised to be there. It is
another to know they are \emph{always} there, for any denominator. Proving that
this decomposition always works is short, concrete, and explains why the cover-up
method succeeds.

\begin{whybox}
Here is the mechanism, in the cleanest case where the denominator
$Q(x) = (x-r_1)(x-r_2)\cdots(x-r_n)$ has $n$ \emph{distinct} real roots and
the numerator $P$ has degree less than $n$. We claim
\[
\frac{P(x)}{Q(x)} = \frac{A_1}{x-r_1} + \frac{A_2}{x-r_2} + \cdots +
\frac{A_n}{x-r_n}
\]
for some constants. Why \emph{must} such constants exist?

It comes down to counting. Clearing denominators turns the claim into
\[
P(x) = A_1\,\frac{Q(x)}{x-r_1} + A_2\,\frac{Q(x)}{x-r_2} + \cdots +
A_n\,\frac{Q(x)}{x-r_n}.
\]
Each $\dfrac{Q(x)}{x-r_k}$ is a polynomial of degree $n-1$. The right side is a
combination of $n$ such polynomials, and as we vary the $A_k$ it ranges over
some subspace of the polynomials of degree $\le n-1$. That space has dimension
$n$ (the coefficients of $1, x, \ldots, x^{n-1}$). We have $n$ knobs $A_k$ and
$n$ coefficients to match -- a square system. So the question is only whether
those $n$ building blocks $\dfrac{Q(x)}{x-r_k}$ are \emph{linearly
independent} -- meaning no block can be formed as a combination of the others,
so each provides a genuinely new direction of freedom; if they are, they span
everything of degree $\le n-1$, and \emph{every} valid $P$ can be hit.

They are independent, and the cover-up trick proves it. Suppose a combination
$\sum_k c_k \dfrac{Q(x)}{x-r_k}$ is the zero polynomial. Evaluate at
$x = r_j$. Every term with $k \neq j$ still contains the factor $(x - r_j)$ and
so vanishes there; the surviving $k=j$ term equals $c_j \prod_{m\neq
j}(r_j - r_m)$, a product of nonzero numbers (the roots are distinct). For the
whole thing to be zero we need $c_j = 0$. This holds for every $j$, so all
coefficients are zero: the building blocks are independent. A square system
with independent columns has exactly one solution for any right-hand side.
\textbf{The decomposition exists, and it is unique.}
\end{whybox}

That argument is worth rereading, because it explains the cover-up method at
the same time as it proves existence: the very evaluation that isolates $A_j$
is the evaluation that shows the building blocks cannot cancel each other.

\begin{connectionbox}
This is a direct application of finite-dimensional linear algebra: $n$ linearly independent
vectors in an $n$-dimensional vector space form a basis and span the entire space. Partial fractions
is one of many places in calculus where a dimension-counting argument guarantees that a desired
representation exists before computing its coefficients -- just as matching $n+1$ derivative conditions
determines a unique Taylor polynomial of degree $n$, or $n+1$ data points determine a unique
interpolating polynomial of degree at most $n$.
\end{connectionbox}

\begin{historybox}
In 1702, writing in \emph{Acta Eruditorum}, Gottfried Wilhelm Leibniz considered
whether every rational function could be integrated into logarithms and arctangents.
He believed he had found a counterexample in the fraction
\[
  \frac{1}{x^4 + 1}.
\]
Factoring over the complex numbers gives $(x^2 + i)(x^2 - i)$. Because $\pm i$ are
imaginary, Leibniz argued that these factors could not be combined into real
quadratic polynomials, concluding that $x^4 + 1$ could not be factored over the
reals and that its integral could not be expressed with real logarithms or arctangents.

Johann Bernoulli pointed out that complex roots of real polynomials always occur in
conjugate pairs, which multiply together to form real quadratics. In 1748, Leonhard
Euler confirmed this by factoring $x^4 + 1$ into real quadratics:
\[
  x^4 + 1 = (x^2 + 1)^2 - 2x^2 = (x^2 + \sqrt{2}x + 1)(x^2 - \sqrt{2}x + 1).
\]
Both quadratic factors are real, and $\int \frac{\dd x}{x^4+1}$ splits into standard
partial fractions. The general principle was established in 1799, when Carl Friedrich
Gauss proved the Fundamental Theorem of Algebra: every real polynomial factors completely
into real linear and real irreducible quadratic factors.
\end{historybox}

\section{The Four Cases of Partial Fractions}
\label{ch24:four-cases}

To decompose \emph{any} proper rational function you handle the denominator's
factors one type at a time. Real-coefficient polynomials factor completely
into linear factors $(x-r)$ and irreducible quadratics $(x^2+px+q)$ (those
with no real root), each possibly repeated. That gives exactly four kinds of
pieces.

\begin{focusbox}
\textbf{The setup rules.} First make the fraction \emph{proper} (numerator
degree $<$ denominator degree) by long division if needed. Then for each factor
of $Q(x)$, add terms:
\begin{itemize}
\item Distinct linear factor $(x-r)$: one term $\dfrac{A}{x-r}$.
\item Repeated linear factor $(x-r)^m$: a full tower
   $\dfrac{A_1}{x-r} + \dfrac{A_2}{(x-r)^2} + \cdots + \dfrac{A_m}{(x-r)^m}$.
\item Distinct irreducible quadratic $(x^2+px+q)$: one term
   $\dfrac{Bx+C}{x^2+px+q}$ (a linear numerator, because the denominator is
   degree two).
\item Repeated irreducible quadratic $(x^2+px+q)^m$: a tower
   $\dfrac{B_1x+C_1}{x^2+px+q} + \cdots + \dfrac{B_mx+C_m}{(x^2+px+q)^m}$.
\end{itemize}
\end{focusbox}

The rule about repeated factors deserves a word of intuition, because students
forget the lower powers.

\begin{intuitionbox}
Why must $(x-r)^3$ contribute three terms, $\dfrac{A_1}{x-r} +
\dfrac{A_2}{(x-r)^2} + \dfrac{A_3}{(x-r)^3}$, and not just the top one? Because
you need enough freedom. A general numerator over $(x-r)^3$ can be \emph{any}
quadratic $a + b(x-r) + c(x-r)^2$ near $x=r$. Dividing that by $(x-r)^3$ gives
exactly $\dfrac{a}{(x-r)^3} + \dfrac{b}{(x-r)^2} + \dfrac{c}{x-r}$. The three
terms are the three coefficients of that local quadratic. Drop any one and you
lose a degree of freedom you might need to match the original numerator.
\end{intuitionbox}

\subsection{Why a quadratic factor gets a linear numerator}

Over an irreducible quadratic the numerator is $Bx + C$, not just a constant.
The reason is the same counting principle: a piece $\dfrac{N(x)}{x^2+px+q}$ is
``proper'' as long as $\deg N < 2$, so the most general allowed numerator has
\emph{two} free constants, $B$ and $C$. Using only a constant would throw away
half the freedom and the system would have no solution in general.

\begin{trapbox}
The three classic errors:
\begin{enumerate}
\item Forgetting to do long division first when the fraction is
   \emph{improper}. $\dfrac{x^3}{x^2-1}$ is not proper; you must divide to get
   $x + \dfrac{x}{x^2-1}$ before decomposing.
\item Writing only $\dfrac{A}{(x-r)^m}$ for a repeated factor and omitting the
   lower powers. The tower must be complete.
\item Putting a constant $A$ over an irreducible quadratic instead of $Bx+C$.
\end{enumerate}
Each of these makes the linear system unsolvable -- a useful early warning:
if your constants come out inconsistent, suspect a malformed setup.
\end{trapbox}

\begin{recoverybox}
\textbf{Pause and predict.} Look at $\dfrac{2x+5}{(x-3)(x^2+1)}$. Before touching any algebra:
how many unknown constants will you need to solve for, and what species of functions will appear
in the final antiderivative?
The denominator degree is $1 + 2 = 3$, so you need three constants ($A$ over $x-3$, and $Bx+C$
over $x^2+1$). The linear factor will produce a logarithm $\ln\abs{x-3}$; the irreducible quadratic
will split into a logarithm $\ln(x^2+1)$ and an arctangent $\arctan x$. Predicting the final shape
before starting keeps you oriented throughout the algebra.
\end{recoverybox}

\subsection{Worked example: a repeated linear factor}

\begin{example}[Repeated root]
\label{ch24:ex-repeated}
Integrate $\displaystyle \int \frac{x + 4}{(x-1)^2(x+2)}\,\dd x.$
\end{example}

The denominator has a repeated factor $(x-1)^2$ and a single factor $(x+2)$,
so the setup is
\[
\frac{x+4}{(x-1)^2(x+2)} = \frac{A}{x-1} + \frac{B}{(x-1)^2} + \frac{D}{x+2}.
\]
Clear denominators:
\[
x + 4 = A(x-1)(x+2) + B(x+2) + D(x-1)^2.
\]
Use the cover-up values for the easy constants. Set $x = 1$: $1 + 4 = B(3)$, so
$B = \tfrac{5}{3}$. Set $x = -2$: $-2 + 4 = D(-3)^2$, so $2 = 9D$ and $D =
\tfrac{2}{9}$. The constant $A$ has no root to cover up (it sits over the
\emph{first} power of a repeated factor), so match one more coefficient. The
$x^2$ coefficient on the left is $0$; on the right it is $A + D$. Hence $A =
-D = -\tfrac{2}{9}$.

So
\[
\frac{x+4}{(x-1)^2(x+2)} = \frac{-2/9}{x-1} + \frac{5/3}{(x-1)^2}
+ \frac{2/9}{x+2},
\]
and integrating term by term,
\[
\int \frac{x+4}{(x-1)^2(x+2)}\,\dd x
= -\frac{2}{9}\ln\abs{x-1} - \frac{5/3}{x-1} + \frac{2}{9}\ln\abs{x+2} + C.
\]
Notice the middle term: the integral of $1/(x-1)^2$ is $-1/(x-1)$, a power, not
a logarithm. Repeated factors produce these non-log pieces, which is exactly
why their tower matters.

\subsection{Worked example: an irreducible quadratic}

\begin{example}[Quadratic factor]
\label{ch24:ex-quadratic}
Integrate $\displaystyle \int \frac{2x^2 - x + 4}{x(x^2+4)}\,\dd x.$
\end{example}

Set up with a linear numerator over the quadratic:
\[
\frac{2x^2 - x + 4}{x(x^2+4)} = \frac{A}{x} + \frac{Bx + C}{x^2 + 4}.
\]
Clear denominators: $2x^2 - x + 4 = A(x^2+4) + (Bx + C)x$. Set $x = 0$:
$4 = 4A$, so $A = 1$. Match $x^2$: $2 = A + B$, so $B = 1$. Match $x^1$:
$-1 = C$. Thus
\[
\frac{2x^2-x+4}{x(x^2+4)} = \frac{1}{x} + \frac{x - 1}{x^2+4}.
\]
Now integrate. The first piece is a logarithm. The second piece we
\emph{split} into a part whose numerator is the derivative of the denominator
and a leftover:
\[
\int \frac{x - 1}{x^2 + 4}\,\dd x
= \int \frac{x}{x^2+4}\,\dd x - \int \frac{1}{x^2+4}\,\dd x.
\]
For the first, $u = x^2 + 4$, $\dd u = 2x\,\dd x$, giving
$\tfrac{1}{2}\ln(x^2+4)$. The second is a standard arctangent with $a = 2$:
$\tfrac{1}{2}\arctan(x/2)$. Altogether,
\[
\int \frac{2x^2-x+4}{x(x^2+4)}\,\dd x
= \ln\abs{x} + \frac{1}{2}\ln(x^2+4) - \frac{1}{2}\arctan\!\frac{x}{2} + C.
\]

\begin{happeningbox}
Look at what the quadratic piece actually became: \emph{a logarithm plus an
arctangent}. This is the universal fate of an irreducible-quadratic term, and
there is a reason. The $Bx$ part of the numerator, after you arrange it to be
a multiple of the denominator's derivative $2x$, integrates to $\ln(x^2+4)$ --
that is substitution. The leftover constant part has no $x$ on top, so it has
nothing to substitute and instead lands on the arctangent form. Every
irreducible quadratic splits its numerator into a ``derivative-of-denominator''
half (a log) and a ``pure constant'' half (an arctangent). Once you see this
split you can integrate any quadratic term on autopilot.
\end{happeningbox}

\subsection{Worked example: improper, needing division first}

\begin{example}[Improper fraction]
\label{ch24:ex-improper}
Integrate $\displaystyle \int \frac{x^3}{x^2 - 1}\,\dd x.$
\end{example}

The numerator degree ($3$) exceeds the denominator degree ($2$), so do
polynomial long division first: $x^3 = x(x^2-1) + x$, hence
\[
\frac{x^3}{x^2-1} = x + \frac{x}{x^2-1}.
\]
The leftover $\dfrac{x}{x^2-1} = \dfrac{x}{(x-1)(x+1)}$ decomposes as
$\dfrac{1/2}{x-1} + \dfrac{1/2}{x+1}$ (cover-up: at $x=1$, $\tfrac{1}{2}$; at
$x=-1$, $\tfrac{-1}{-2} = \tfrac{1}{2}$). So
\[
\int \frac{x^3}{x^2-1}\,\dd x
= \frac{x^2}{2} + \frac{1}{2}\ln\abs{x-1} + \frac{1}{2}\ln\abs{x+1} + C
= \frac{x^2}{2} + \frac{1}{2}\ln\abs{x^2-1} + C.
\]

\begin{walkbox}
On your walk, picture why long division is unavoidable. A proper fraction
$P/Q$ with $\deg P < \deg Q$ shrinks toward $0$ as $x\to\infty$. But $x^3/(x^2-1)
\approx x$ for large $x$ -- it grows without bound. No sum of pieces like
$\dfrac{A}{x-r}$ (each shrinking to zero) could ever reproduce that growth. The
division peels off the polynomial part $x$ that carries the growth, leaving a
genuinely proper remainder for the fractions to handle. The methods only work
on the part of the function that actually decays.
\end{walkbox}

\subsection{The shape of the picture}

It helps to \emph{see} a decomposition. Below, the rational function
$\dfrac{3x+5}{(x-2)(x+1)}$ from earlier is drawn together with its two simple
pieces. Near each vertical asymptote, one piece dominates and the other is
nearly flat: near $x=2$ the curve behaves like $\tfrac{11/3}{x-2}$, and near $x=-1$
it behaves like $-\tfrac{2/3}{x+1}$.

\begin{figure}[H]
\centering
\begin{tikzpicture}
\begin{axis}[
  width=12cm, height=7cm,
  axis lines=middle,
  xlabel=$x$, ylabel=$y$,
  xmin=-4, xmax=5, ymin=-8, ymax=8,
  restrict y to domain=-30:30,
  samples=200,
  xtick={-4,-2,1,2,3,4,5},
  xticklabel style={fill=white, inner sep=1pt, font=\footnotesize},
  yticklabel style={fill=white, inner sep=1pt, font=\footnotesize},
  axis on top,
  legend pos=north east, legend cell align=left,
  every axis plot/.append style={very thick},
]
\addplot[curveblue, domain=-4:-1.1] {(3*x+5)/((x-2)*(x+1))};
\addlegendentry{$\frac{3x+5}{(x-2)(x+1)}$}
\addplot[curveblue, forget plot, domain=-0.9:1.9] {(3*x+5)/((x-2)*(x+1))};
\addplot[curveblue, forget plot, domain=2.1:5] {(3*x+5)/((x-2)*(x+1))};
\addplot[tangentred, dashed, domain=-4:1.9] {(11/3)/(x-2)};
\addlegendentry{$\frac{11/3}{x-2}$}
\addplot[tangentred, dashed, forget plot, domain=2.1:5] {(11/3)/(x-2)};
\addplot[targetgreen, dotted, domain=-4:-1.1] {(-2/3)/(x+1)};
\addlegendentry{$-\frac{2/3}{x+1}$}
\addplot[targetgreen, dotted, forget plot, domain=-0.9:5] {(-2/3)/(x+1)};
\end{axis}
\end{tikzpicture}
\caption{A rational function (solid) and its two partial-fraction pieces
(dashed, dotted). Each piece controls the behavior near its own asymptote; the
sum reproduces the original everywhere.}
\label{ch24:fig-decomp}
\end{figure}

\section{Trigonometric Integrals: Powers and Products}
\label{ch24:trig-integrals}

We turn now to integrals built from powers and products of sines and cosines, such as
$\int \sin^5 x\,\dd x$, $\int \sin^2 x\cos^2 x\,\dd x$, or
$\int \sin 3x\cos 5x\,\dd x$. None is a standard form as written. But every one
yields to the same move: \emph{use a trig identity to reshape the
integrand into something a substitution can finish}. There is no new calculus,
only the Pythagorean identity, the double-angle formulas, and the
product-to-sum formulas, deployed cleverly.

\subsection{Odd powers: peel off one factor}

\begin{example}[Odd power of sine]
\label{ch24:ex-odd}
Integrate $\displaystyle \int \sin^5 x\,\dd x.$
\end{example}

There is an \emph{odd} number of sines, so peel one off
and convert the rest to cosines using $\sin^2 x = 1 - \cos^2 x$. The lone sine
becomes the $\dd u$ for $u = \cos x$.
\begin{align*}
\int \sin^5 x\,\dd x
&= \int (\sin^2 x)^2 \sin x\,\dd x
= \int (1 - \cos^2 x)^2 \sin x\,\dd x.
\end{align*}
Substitute $u = \cos x$, $\dd u = -\sin x\,\dd x$:
\[
= -\int (1 - u^2)^2\,\dd u
= -\int (1 - 2u^2 + u^4)\,\dd u
= -\Big(u - \tfrac{2}{3}u^3 + \tfrac{1}{5}u^5\Big) + C.
\]
Back-substitute $u = \cos x$:
\[
\int \sin^5 x\,\dd x = -\cos x + \frac{2}{3}\cos^3 x - \frac{1}{5}\cos^5 x + C.
\]

\begin{whybox}
Why does peeling one factor work, and why does it need an \emph{odd} power?
Substitution $u=\cos x$ needs $\dd u = -\sin x\,\dd x$ -- it consumes exactly
\emph{one} sine. After spending that one, the remaining sines must be
expressible in $u=\cos x$, and $\sin^2 x = 1 - \cos^2 x$ does that perfectly
because it only ever uses sine in \emph{even} amounts. So you need one sine to
spend on $\dd u$ and an even number left over to convert: total odd. If the
power were even, you would have a leftover lone sine with nothing to pair it
with, and the substitution stalls. Odd powers are easy precisely because they
split as ``one for $\dd u$, the rest even.''
\end{whybox}

The same trick handles odd powers of cosine (peel one cosine, set $u=\sin x$)
and mixed products $\int \sin^m x\cos^n x\,\dd x$ \emph{whenever at least one
of $m,n$ is odd}: peel one factor off the odd one, convert the rest, substitute
for the other function.

\begin{example}[Mixed product, one odd power]
$\displaystyle \int \sin^2 x\cos^3 x\,\dd x.$ The cosine power is odd. Peel one
cosine, convert the rest with $\cos^2 x = 1 - \sin^2 x$, set $u = \sin x$:
\[
\int \sin^2 x(1-\sin^2 x)\cos x\,\dd x
= \int u^2(1-u^2)\,\dd u
= \frac{u^3}{3} - \frac{u^5}{5}
= \frac{\sin^3 x}{3} - \frac{\sin^5 x}{5} + C.
\]
\end{example}

\begin{recoverybox}
\textbf{Pause and predict.} What is the definite integral $\int_0^{2\pi} \sin^7 x\,\dd x$?
Before reaching for pencil and paper: sine oscillates symmetrically above and below
zero. Every positive peak on $[0, \pi]$ is matched by an identical negative
trough on $[\pi, 2\pi]$. Raising sine to an odd power preserves this sign
reversal ($(-y)^7 = -y^7$), so the two halves cancel completely: the integral is $0$
by pure symmetry.
\end{recoverybox}

\subsection{Even powers: the half-angle slowdown}

When \emph{both} powers are even, the peel-one trick fails -- there is no spare
factor to be a $\dd u$. Instead we lower the powers using the half-angle
identities, which are the double-angle formula solved for the squares:
\[
\sin^2 x = \frac{1 - \cos 2x}{2},
\qquad
\cos^2 x = \frac{1 + \cos 2x}{2}.
\]
Each application trades a squared sine or cosine for a plain cosine of double
the angle -- lower power, in exchange for a faster oscillation. Repeat until
nothing is squared.

\begin{example}[Even power]
\label{ch24:ex-even}
Integrate $\displaystyle \int \sin^2 x\,\dd x.$
\[
\int \sin^2 x\,\dd x
= \int \frac{1 - \cos 2x}{2}\,\dd x
= \frac{1}{2}\Big(x - \frac{\sin 2x}{2}\Big) + C
= \frac{x}{2} - \frac{\sin 2x}{4} + C.
\]
\end{example}

\begin{intuitionbox}
There is a sanity check hiding here. The average value of $\sin^2 x$ over a
full period is $\tfrac{1}{2}$ -- because $\sin^2 + \cos^2 = 1$ and, by
symmetry, $\sin^2$ and $\cos^2$ have the same average, so each must average
$\tfrac{1}{2}$. Our antiderivative $\tfrac{x}{2} - \tfrac{\sin 2x}{4}$ grows
like $\tfrac{x}{2}$ with a small wobble on top, so over a period it
accumulates exactly $\tfrac{1}{2}\cdot(\text{length})$. The formula encodes
the average. This ``$\sin^2$ averages to a half'' fact is the reason the RMS
of an alternating current is the peak over $\sqrt 2$.
\end{intuitionbox}

\begin{example}[Both even]
Integrate $\displaystyle \int \sin^2 x\cos^2 x\,\dd x.$ The slick route uses the double-angle identity
$\sin x\cos x = \tfrac{1}{2}\sin 2x$, so $\sin^2 x\cos^2 x =
\tfrac{1}{4}\sin^2 2x$. Then half-angle again on $\sin^2 2x$:
\[
\int \tfrac{1}{4}\sin^2 2x\,\dd x
= \frac{1}{4}\int \frac{1 - \cos 4x}{2}\,\dd x
= \frac{x}{8} - \frac{\sin 4x}{32} + C.
\]
\end{example}

\subsection{Products of different frequencies}

For $\int \sin(mx)\cos(nx)\,\dd x$ with $m \neq n$ there are no powers at all,
just a product of two waves at different frequencies. The reshaping tool is the
\emph{product-to-sum} identities:
\[
\sin A\cos B = \tfrac{1}{2}\big[\sin(A+B) + \sin(A-B)\big],
\]
and its siblings for $\cos A\cos B$ and $\sin A\sin B$. A product of two waves
becomes a \emph{sum} of two waves, and a sum is trivial to integrate.

\begin{example}
$\displaystyle \int \sin 3x\cos 5x\,\dd x
= \frac{1}{2}\int \big[\sin 8x + \sin(-2x)\big]\,\dd x
= -\frac{\cos 8x}{16} + \frac{\cos 2x}{4} + C.$
\end{example}

\begin{connectionbox}
These product-to-sum integrals are the computational core of \textbf{Fourier
series} -- the decomposition of an arbitrary wave into a sum of pure harmonic
sines and cosines. The reason you can extract the ``amount of frequency $n$''
inside a signal is that $\int_0^{2\pi}\sin(mx)\sin(nx)\,\dd x = 0$ whenever
$m \neq n$ and equals $\pi$ when $m = n$. That \emph{orthogonality} (two
functions are called orthogonal when their product integrates to zero over a
full period, exactly like two perpendicular geometric vectors whose dot product
is zero) is the direct result of the product-to-sum computation above:
different frequencies multiply to a sum of waves that integrates to zero over
a cycle; the same frequency multiplies to a squared wave with a nonzero
average $\tfrac12$. Every time you stream music, compress audio, or filter
signals, this integral is doing the work.
\end{connectionbox}

\begin{figure}[H]
\centering
\begin{tikzpicture}
\begin{axis}[
  width=12cm, height=5.5cm,
  axis lines=middle,
  xlabel=$x$, ylabel=$y$,
  xmin=0, xmax=6.8, ymin=-0.25, ymax=1.35,
  domain=0:6.283, samples=200,
  xtick={0,3.1416,6.283}, xticklabels={$0$,$\pi$,$2\pi$},
  ytick={0,0.5,1}, yticklabels={$0$,$\tfrac12$,$1$},
  every axis plot/.append style={very thick},
]
\addplot[curveblue] {sin(deg(2*x))*sin(deg(2*x))};
\node[curveblue, anchor=south, font=\small] at (axis cs:2.35, 1.03) {$y=\sin^2 2x$};
\addplot[tangentred, dashed] {0.5};
\node[tangentred, anchor=south east, font=\small] at (axis cs:6.7, 0.53) {$\text{average}=\tfrac12$};
\end{axis}
\end{tikzpicture}
\caption{A squared sine oscillates between $0$ and $1$ but spends its time
symmetrically, averaging $\tfrac12$. The half-angle identity
$\sin^2 2x = \tfrac{1-\cos 4x}{2}$ makes this visible: a constant $\tfrac12$
plus a zero-average wobble.}
\label{ch24:fig-sin2}
\end{figure}

\section{Completing the Square: Taming Any Quadratic}
\label{ch24:complete-square}

Before trig substitution we need one small, powerful algebraic move. Many
integrals hide a quadratic in the denominator or under a root that is not
quite a standard form, like
\[
\int \frac{\dd x}{x^2 + 4x + 13}.
\]
The denominator is not $x^2 + a^2$, so the arctangent form does not apply
directly. But \emph{completing the square} reshapes any quadratic into a
perfect square plus (or minus) a constant:
\[
x^2 + 4x + 13 = (x^2 + 4x + 4) + 9 = (x+2)^2 + 9 = (x+2)^2 + 3^2.
\]
Now substitute $u = x + 2$ (a pure shift, $\dd u = \dd x$) and it \emph{is} the
arctangent form with $a = 3$:
\[
\int \frac{\dd x}{(x+2)^2 + 3^2}
= \frac{1}{3}\arctan\!\frac{x+2}{3} + C.
\]

\begin{happeningbox}
Completing the square amounts geometrically to a \emph{horizontal shift}. The graph of
$y = x^2 + 4x + 13$ is the parabola $y = x^2 + 9$ shifted two units to the left
(Figure~\ref{ch24:fig-complete-square}). The substitution $u = x+2$ slides
our coordinate system to sit at the bottom of the parabola, where the quadratic
is symmetric: $u^2 + 9$. Every quadratic is a standard parabola shifted horizontally
and vertically; completing the square determines the shift, and the substitution
undoes it. Once centered, the integrand matches the standard arctangent form.
\end{happeningbox}

\begin{figure}[H]
\centering
\begin{tikzpicture}
\begin{axis}[
  width=11cm, height=5.8cm,
  axis lines=middle,
  xlabel=$x$, ylabel=$y$,
  xmin=-5.5, xmax=3.5, ymin=0, ymax=24,
  domain=-5:3, samples=120,
  xtick={-2,0}, xticklabels={$-2$,$0$},
  ytick=\empty,
  legend style={at={(0.5,1.05)}, anchor=south, legend columns=2, font=\footnotesize, /tikz/every even column/.append style={column sep=12pt}},
  legend cell align=left,
  every axis plot/.append style={very thick},
]
\addplot[tangentred, domain=-5:1] {(x+2)^2 + 9};
\addlegendentry{$y = (x+2)^2 + 9 = x^2+4x+13$}
\addplot[curveblue, dashed, domain=-3:3] {x^2 + 9};
\addlegendentry{$y = x^2 + 9$}
\draw[dotted, thick, gray] (axis cs:-2,0) -- (axis cs:-2,9);
\draw[dotted, thick, gray] (axis cs:0,0) -- (axis cs:0,9);
\draw[-{Stealth[length=5pt]}, thick, quotientorange] (axis cs:-1.9,4.2) -- node[midway, above, font=\footnotesize, text=quotientorange] {shift $u = x+2$} (axis cs:-0.1,4.2);
\node[circle, fill=tangentred, inner sep=1.5pt] at (axis cs:-2,9) {};
\node[circle, fill=curveblue, inner sep=1.5pt] at (axis cs:0,9) {};
\node[anchor=south east, font=\footnotesize] at (axis cs:-0.08,9) {$9$};
\end{axis}
\end{tikzpicture}
\caption{Completing the square reveals the hidden symmetry of a quadratic. The parabola
$y = x^2+4x+13$ (solid) is identical to the centered parabola $y = x^2+9$ (dashed), shifted
two units to the left. The linear substitution $u = x+2$ recenters the origin at the vertex,
turning the un-factored quadratic into the standard sum of squares $u^2+3^2$.}
\label{ch24:fig-complete-square}
\end{figure}

\begin{example}[Linear numerator over a completed square]
\label{ch24:ex-complete}
$\displaystyle \int \frac{x + 1}{x^2 + 4x + 13}\,\dd x.$ Split the numerator into
a multiple of the denominator's derivative $2x+4$ plus a leftover constant:
$x + 1 = \tfrac{1}{2}(2x+4) - 1$. Then
\[
\int \frac{x+1}{x^2+4x+13}\,\dd x
= \frac{1}{2}\int \frac{2x+4}{x^2+4x+13}\,\dd x - \int \frac{\dd x}{x^2+4x+13}.
\]
The first integral is $\tfrac{1}{2}\ln(x^2+4x+13)$ by substitution; the second
is the arctangent we just did. So the answer is
\[
\frac{1}{2}\ln(x^2+4x+13) - \frac{1}{3}\arctan\!\frac{x+2}{3} + C.
\]
\end{example}

This ``log half plus arctangent half'' split is the same one from the partial
fractions section -- another sign that all these techniques are one idea.

\section{Trigonometric Substitution: Algebra Becomes Geometry}
\label{ch24:trig-sub}

The last technique is the most geometric one in the chapter. Integrals
containing $\sqrt{a^2 - x^2}$, $\sqrt{a^2 + x^2}$, or $\sqrt{x^2 - a^2}$ resist
ordinary substitution because the square root has no convenient derivative
hiding in the integrand. The escape is to \emph{change variables so that the
expression under the root becomes a perfect square}, using the Pythagorean
identities. And the cleanest way to see which substitution to use, and how to
get back, is to \emph{draw a right triangle}.

\subsection{The idea on a triangle}

Consider $\sqrt{a^2 - x^2}$. The relation $a^2 = x^2 + (\sqrt{a^2-x^2})^2$ is
the Pythagorean theorem: $a$ is a hypotenuse, $x$ and $\sqrt{a^2-x^2}$ are the
two legs. So \emph{draw} that triangle and let $\theta$ be the angle whose
\emph{opposite leg} is $x$:
\[
\sin\theta = \frac{x}{a} \;\Longrightarrow\; x = a\sin\theta,
\qquad
\sqrt{a^2 - x^2} = a\cos\theta.
\]
To make the substitution reversible, we restrict $\theta$ to the principal
interval $[-\pi/2, \pi/2]$. On this interval, cosine is always nonnegative, so
$\sqrt{\cos^2\theta} = \abs{\cos\theta} = \cos\theta$ -- no absolute-value sign
needed.

\begin{figure}[H]
\centering
\begin{tikzpicture}[scale=0.9]
% Triangle 1: x = a sin theta
\begin{scope}[shift={(0,0)}]
  \coordinate (O1) at (0,0);
  \coordinate (A1) at (3.4,0);
  \coordinate (B1) at (3.4,2.2);
  \draw[very thick, sourceblue] (O1) -- (A1) -- (B1) -- cycle;
  \draw (A1) ++(-0.25,0) -- ++(0,0.25) -- ++(0.25,0);
  \node[below, sourceblue, font=\small] at (1.7,-0.1) {$\sqrt{a^2-x^2} = a\cos\theta$};
  \node[right, tangentred, font=\small] at (3.4,1.1) {$x = a\sin\theta$};
  \node[above left, structurepurple, font=\small] at (1.6,1.2) {$a$};
  \draw (O1)++(0.6,0) arc (0:33:0.6);
  \node[font=\footnotesize] at (0.85,0.2) {$\theta$};
  \node[above, font=\footnotesize\bfseries, align=center] at (1.7,2.6)
    {\textbf{(a)} $\sqrt{a^2 - x^2}$\\$x = a\sin\theta$};
\end{scope}

% Triangle 2: x = a tan theta
\begin{scope}[shift={(5.4,0)}]
  \coordinate (O2) at (0,0);
  \coordinate (A2) at (2.8,0);
  \coordinate (B2) at (2.8,2.2);
  \draw[very thick, targetgreen!80!black] (O2) -- (A2) -- (B2) -- cycle;
  \draw (A2) ++(-0.25,0) -- ++(0,0.25) -- ++(0.25,0);
  \node[below, sourceblue, font=\small] at (1.4,-0.1) {$a$};
  \node[right, tangentred, font=\small] at (2.8,1.1) {$x = a\tan\theta$};
  \node[above left, quotientorange!90!black, font=\small] at (1.3,1.2) {$\sqrt{a^2+x^2} = a\sec\theta$};
  \draw (O2)++(0.6,0) arc (0:38:0.6);
  \node[font=\footnotesize] at (0.85,0.22) {$\theta$};
  \node[above, font=\footnotesize\bfseries, align=center] at (1.4,2.6)
    {\textbf{(b)} $\sqrt{a^2 + x^2}$\\$x = a\tan\theta$};
\end{scope}

% Triangle 3: x = a sec theta
\begin{scope}[shift={(10.7,0)}]
  \coordinate (O3) at (0,0);
  \coordinate (A3) at (2.4,0);
  \coordinate (B3) at (2.4,2.2);
  \draw[very thick, quotientorange] (O3) -- (A3) -- (B3) -- cycle;
  \draw (A3) ++(-0.25,0) -- ++(0,0.25) -- ++(0.25,0);
  \node[below, sourceblue, font=\small] at (1.2,-0.1) {$a$};
  \node[right, targetgreen!80!black, font=\small] at (2.4,1.1) {$\sqrt{x^2-a^2} = a\tan\theta$};
  \node[above left, structurepurple, font=\small] at (1.1,1.2) {$x = a\sec\theta$};
  \draw (O3)++(0.5,0) arc (0:43:0.5);
  \node[font=\footnotesize] at (0.75,0.25) {$\theta$};
  \node[above, font=\footnotesize\bfseries, align=center] at (1.2,2.6)
    {\textbf{(c)} $\sqrt{x^2 - a^2}$\\$x = a\sec\theta$};
\end{scope}

\end{tikzpicture}
\caption{The three canonical right-triangle substitutions. Each substitution is engineered so the Pythagorean theorem collapses the square root to a single trigonometric function: \textbf{(a)} $x=a\sin\theta$ yields $a^2-x^2 = a^2\cos^2\theta$; \textbf{(b)} $x=a\tan\theta$ yields $a^2+x^2 = a^2\sec^2\theta$; \textbf{(c)} $x=a\sec\theta$ yields $x^2-a^2 = a^2\tan^2\theta$. Reading the triangle's ratios translates any trig expression of $\theta$ back to $x$.}
\label{ch24:fig-triangle}
\end{figure}

The reference triangle tracks the relationships among the sides. It tells
you the substitution ($x = a\sin\theta$), it tells you $\dd x = a\cos\theta\,
\dd\theta$, it converts the root to $a\cos\theta$, and most importantly, it tells
you how to convert the \emph{answer} back to $x$ at the end, because every trig
function of $\theta$ is just a ratio of the triangle's sides.

\begin{focusbox}
\textbf{The three triangles.}
\begin{itemize}
\item $\sqrt{a^2 - x^2}$: set $x = a\sin\theta$, then root $= a\cos\theta$.
   ($\sin^2 + \cos^2 = 1$, rearranged.)
\item $\sqrt{a^2 + x^2}$: set $x = a\tan\theta$, then root $= a\sec\theta$.
   ($1 + \tan^2 = \sec^2$.)
\item $\sqrt{x^2 - a^2}$: set $x = a\sec\theta$, then root $= a\tan\theta$.
   ($\sec^2 - 1 = \tan^2$.)
\end{itemize}
In each case the substitution is the one that makes a Pythagorean identity
collapse the root into a single trig function. Draw the triangle and you never
have to memorize which is which.
\end{focusbox}

\begin{whybox}
Why does turning $x$ into a trig function help at all? Because the obstacle is
the square root, and a square root vanishes only when its inside is a perfect
square. The Pythagorean identities are precisely the algebraic statements that
$a^2 - (a\sin\theta)^2$, $a^2 + (a\tan\theta)^2$, and $(a\sec\theta)^2 - a^2$
are each a perfect square ($a^2\cos^2\theta$, $a^2\sec^2\theta$,
$a^2\tan^2\theta$). Trig substitution is the deliberate choice of a new
variable engineered so that the thing under the root becomes a square and the
root opens up cleanly. It is substitution aimed by geometry.
\end{whybox}

\subsection{Worked example: the area of a circle, recovered}

\begin{example}[Root of $a^2 - x^2$]
\label{ch24:ex-circle}
Integrate $\displaystyle \int \sqrt{a^2 - x^2}\,\dd x.$
\end{example}

This integrand is the height of a circle of radius $a$, so the integral is a
circular area (Figure~\ref{ch24:fig-semicircle}); watch the geometry fall out
of the algebra. Use the triangle: $x = a\sin\theta$, $\dd x =
a\cos\theta\,\dd\theta$, and $\sqrt{a^2-x^2} = a\cos\theta$. The integral
becomes pure trig:
\[
\int a\cos\theta\cdot a\cos\theta\,\dd\theta
= a^2\int \cos^2\theta\,\dd\theta
= a^2\Big(\frac{\theta}{2} + \frac{\sin 2\theta}{4}\Big) + C,
\]
using the even-power technique from before. Now convert back via the triangle.
From $\sin\theta = x/a$ we get $\theta = \arcsin(x/a)$, and $\sin 2\theta =
2\sin\theta\cos\theta = 2\cdot\dfrac{x}{a}\cdot\dfrac{\sqrt{a^2-x^2}}{a} =
\dfrac{2x\sqrt{a^2-x^2}}{a^2}$. Therefore
\[
\int \sqrt{a^2-x^2}\,\dd x
= \frac{a^2}{2}\arcsin\!\frac{x}{a} + \frac{x}{2}\sqrt{a^2 - x^2} + C.
\]

\begin{aestheticbox}
Test this on a whole circle. The area inside $x^2 + y^2 = a^2$ is four times
the area under $y = \sqrt{a^2 - x^2}$ from $0$ to $a$:
\[
4\int_0^a \sqrt{a^2 - x^2}\,\dd x
= 4\left[\frac{a^2}{2}\arcsin\!\frac{x}{a} + \frac{x}{2}\sqrt{a^2-x^2}
\right]_0^a.
\]
At $x = a$: $\arcsin 1 = \tfrac{\pi}{2}$ and the square-root term is $0$. At
$x = 0$: everything is $0$. So the area is $4\cdot\dfrac{a^2}{2}\cdot
\dfrac{\pi}{2} = \pi a^2$. The most famous formula in geometry falls out of a
triangle substitution. The $\arcsin$ that appears is
\emph{the angle}, and integrating around the circle is literally accumulating
that angle.
\end{aestheticbox}

\begin{figure}[H]
\centering
\begin{tikzpicture}
\begin{axis}[
  width=11cm, height=5.5cm,
  axis lines=middle,
  xlabel=$x$, ylabel=$y$,
  xmin=-2.4, xmax=2.4, ymin=-0.3, ymax=2.6,
  domain=-2:2, samples=200,
  xtick={-2,-1,0,1,2}, xticklabels={$-a$,,$0$,,$a$},
  ytick={2}, yticklabels={$a$},
  yticklabel style={anchor=south east, font=\footnotesize},
  every axis plot/.append style={very thick},
]
\addplot[curveblue, name path=arc] {sqrt(4 - x*x)};
\path[name path=axisline] (axis cs:-2,0) -- (axis cs:2,0);
\addplot[areafill!70, opacity=0.55] fill between[of=arc and axisline];
\end{axis}
\end{tikzpicture}
\caption{The integral $\int_{-a}^{a}\sqrt{a^2-x^2}\,\dd x$ is the shaded area
under the upper semicircle of radius $a$ (here $a=2$). The trig-substitution
antiderivative evaluates it to $\tfrac{\pi a^2}{2}$, the half-disc area; double
it (or take four quarters, as in the box above) and the full $\pi a^2$ appears.}
\label{ch24:fig-semicircle}
\end{figure}

\subsection{Worked example: a secant integral via the tangent triangle}

\begin{example}[Root of $a^2 + x^2$]
\label{ch24:ex-tan}
Integrate $\displaystyle \int \frac{\dd x}{\sqrt{x^2 + 9}}.$
\end{example}

Here $a = 3$, root is $\sqrt{a^2 + x^2}$, so the tangent triangle:
$x = 3\tan\theta$, $\dd x = 3\sec^2\theta\,\dd\theta$, and $\sqrt{x^2+9} =
3\sec\theta$. Then
\[
\int \frac{3\sec^2\theta}{3\sec\theta}\,\dd\theta
= \int \sec\theta\,\dd\theta
= \ln\abs{\sec\theta + \tan\theta} + C.
\]
Convert back with the triangle. With opposite $= x$, adjacent $= 3$,
hypotenuse $= \sqrt{x^2+9}$, we have $\tan\theta = x/3$ and $\sec\theta =
\sqrt{x^2+9}/3$. So
\[
\int \frac{\dd x}{\sqrt{x^2+9}}
= \ln\left|\frac{\sqrt{x^2+9}}{3} + \frac{x}{3}\right| + C
= \ln\big|x + \sqrt{x^2+9}\big| + C',
\]
absorbing the constant $-\ln 3$ into $C'$. (This is also $\operatorname{arcsinh}(x/3)$
up to a constant -- the inverse hyperbolic sine -- which is no coincidence:
hyperbolic functions parametrize $\sqrt{x^2+a^2}$ the way circular ones
parametrize $\sqrt{a^2-x^2}$.)

\begin{trapbox}
The single most forgotten step in trig substitution is \textbf{converting back
to $x$}. Your answer must be a function of the original variable; an answer
left in $\theta$ is unfinished. Always draw the triangle, label its three
sides from the substitution, and read off whatever trig functions appear in
your antiderivative as ratios of those sides. The triangle is your dictionary
back to $x$.

A second trap: when the integral is \emph{definite}, you may instead change the
limits to $\theta$-values and never convert back -- often cleaner. Pick one
discipline and follow it through.
\end{trapbox}

\subsection{When the root hides inside a shifted quadratic}

Completing the square and trig substitution combine to handle the general
quadratic-under-a-root. For $\int \dfrac{\dd x}{\sqrt{x^2 + 4x + 13}}$, first
complete the square: $x^2 + 4x + 13 = (x+2)^2 + 9$. Substitute $u = x+2$ to
center it, giving $\int \dfrac{\dd u}{\sqrt{u^2 + 9}}$, which is exactly the
tangent-triangle integral above. Reshape (complete the square), then reshape
again (trig substitution): the chapter's whole theme in one problem.

\begin{labbox}
Do this one by hand to feel the pipeline. Integrate
$\displaystyle \int \frac{\dd x}{\sqrt{5 - 4x - x^2}}.$
\begin{enumerate}
\item Complete the square inside the root. Write $5 - 4x - x^2 =
   -(x^2 + 4x) + 5 = -\big((x+2)^2 - 4\big) + 5 = 9 - (x+2)^2$.
\item Substitute $u = x+2$: the integral is $\displaystyle\int
   \frac{\dd u}{\sqrt{9 - u^2}}$, a standard arcsine with $a = 3$.
\item Read off $\arcsin(u/3) = \arcsin\!\dfrac{x+2}{3} + C$.
\end{enumerate}
Notice you did not even need the triangle here -- once centered, it was a
named standard form. Recognizing the standard form is always the fastest path;
the triangle is the fallback when no name fits.
\end{labbox}

\section{One Worldview, Several Tools}
\label{ch24:synthesis}

All of the techniques in this chapter share a single strategy.

\begin{focusbox}
\textbf{Every technique reshapes the integrand into a known form.}
\begin{itemize}
\item \emph{Partial fractions}: a ratio of polynomials $\to$ a sum of
   $\dfrac{A}{(x-r)^k}$ and $\dfrac{Bx+C}{(x^2+px+q)^k}$, each a log, a power,
   or an arctangent.
\item \emph{Trig powers (odd)}: peel a factor $\to$ a polynomial in $\sin$ or
   $\cos$, finished by substitution.
\item \emph{Trig powers (even)}: half-angle identities $\to$ lower powers and
   doubled angles, repeated to the end.
\item \emph{Products of waves}: product-to-sum $\to$ a sum of single waves.
\item \emph{Completing the square}: a general quadratic $\to$ a centered
   $u^2 \pm a^2$, matching the arctangent or arcsine form.
\item \emph{Trig substitution}: a square root $\to$ a triangle, where the
   Pythagorean identity makes the root a perfect square.
\end{itemize}
The calculus underneath never changes. What changes is your fluency at
recognizing which dictionary to open.
\end{focusbox}

\begin{connectionbox}
By the middle of the eighteenth century, Leonhard Euler had organized essentially all of
these techniques into a coherent program. His \emph{Institutiones Calculi
Integralis} (1768) treats rational functions by partial fractions, trig
integrals by the identities, and irrational integrands by substitutions that
rationalize them -- including the trigonometric substitutions for roots of quadratics.
The central idea Euler systematized is the one we have pursued: a wide variety of
integrals can be evaluated not by inventing new operations, but by \emph{transforming integrands
into a small canon of solvable standard forms}. The skill of integration consists largely in
recognizing which algebraic transformation exposes the standard form.
\end{connectionbox}

\begin{aestheticbox}
Notice how small a set of functions suffices here: integrating rational
functions never requires anything beyond rational terms, logarithms, and arctangents,
while integrating square roots of quadratics adds at most arcsines
and square roots. In every case, the antiderivative stays within the class of elementary functions
(combinations of powers, exponentials, logarithms, and trigonometric functions under arithmetic
and composition). By contrast, Liouville proved in 1835 that simple integrands such as
$\int e^{-x^2}\,\dd x$ have no elementary antiderivative at all, whereas every rational function
and every square root of a quadratic integrates systematically within this compact family.
\end{aestheticbox}

\section{Limitations and next steps}
\label{ch24:road-ahead}

\begin{connectionbox}
Partial fractions provides an algorithmic method for integrating any rational function
$P(x)/Q(x)$, provided the denominator $Q(x)$ can be factored into linear and quadratic terms.
Coupled with trigonometric substitution, it also handles many integrands containing
square roots of quadratics.

The method has two main limitations:
\begin{enumerate}
  \item \emph{Factoring high-degree denominators:} Factoring $Q(x)$ can be difficult or
  algebraically impossible. The Abel--Ruffini theorem guarantees that for general polynomials
  of degree $5$ or higher, roots cannot be written in radicals. When factors cannot
  be found explicitly, partial fractions cannot be set up.
  \item \emph{Non-elementary antiderivatives:} Many common integrands, such as $e^{-x^2}$
  or $\sqrt{1+x^4}$, do not possess elementary antiderivatives in terms of standard
  polynomial, exponential, or trigonometric functions.
\end{enumerate}

When an exact symbolic antiderivative is unavailable or unneeded, we turn to
\emph{numerical integration} (Chapter~27), which computes definite integrals directly
from sample points. Before that, Chapter~25 extends integration to unbounded intervals
and vertical asymptotes (\emph{improper integrals}), and Chapter~26 applies integration
to geometric and physical quantities.
\end{connectionbox}

\section{Exercises}
\label{ch24:exercises}

\begin{exercise}[Mental check: Partial fraction setup and counting]
Without solving for any constants, write down the partial fraction decomposition template
for $\dfrac{3x^2 + 5}{(x-2)^2(x^2+9)}$. State how many total constants you must determine,
and predict which pieces will yield logarithms, which will yield powers, and which will
yield arctangents.
\end{exercise}

\begin{exercise}[Mental check: Symmetry and parity in trig integrals]
Without computing any antiderivatives, determine the exact value of
$\displaystyle \int_{-\pi/2}^{\pi/2} \sin^5 x \cos^4 x\,\dd x$, and explain why using the concept of odd functions.
Then predict the exact value of $\displaystyle \int_0^{\pi/2} \sin^2 x\,\dd x$ using the average value of $\sin^2 x$.
\end{exercise}

\begin{exercise}[Distinct linear factors]
Decompose and integrate $\displaystyle \int \frac{5x - 4}{(x-1)(x+3)}\,\dd x.$
Use the cover-up method to find both constants, and verify your decomposition
by recombining over a common denominator.
\end{exercise}

\begin{exercise}[Repeated factor]
Integrate $\displaystyle \int \frac{2x^2 + 3}{(x+1)^3}\,\dd x.$ Set up the full
tower of three terms, find the constants (try the substitution $u = x+1$ as an
alternative to matching coefficients), and identify which pieces become
logarithms and which become powers.
\end{exercise}

\begin{exercise}[Irreducible quadratic]
Integrate $\displaystyle \int \frac{x^2 + x + 1}{x(x^2+1)}\,\dd x.$ Decompose
with a linear numerator over the quadratic, then split that piece into its
``log half'' and ``arctangent half'' as in
Example~\ref{ch24:ex-quadratic}.
\end{exercise}

\begin{exercise}[Improper fraction]
Integrate $\displaystyle \int \frac{x^2}{x^2 - 4}\,\dd x.$ First explain, in
one sentence, why long division is required before any decomposition, then
carry it out.
\end{exercise}

\begin{exercise}[Odd trig power]
Integrate $\displaystyle \int \cos^3 x\,\dd x$ and
$\displaystyle \int \sin^3 x\cos^4 x\,\dd x.$ State which factor you peel and
why the substitution then closes.
\end{exercise}

\begin{exercise}[Even trig powers]
Integrate $\displaystyle \int \cos^4 x\,\dd x$ by applying the half-angle
identity twice. Then, without computing, predict the average value of
$\cos^4 x$ over a full period and check that your antiderivative's linear part
agrees with it.
\end{exercise}

\begin{exercise}[Product of waves]
Integrate $\displaystyle \int \cos 2x\cos 6x\,\dd x$ using a product-to-sum
identity. Then explain, using your answer, why $\int_0^{2\pi}\cos 2x\cos 6x\,
\dd x = 0$, and connect this to the orthogonality of distinct frequencies.
\end{exercise}

\begin{exercise}[Completing the square]
Integrate $\displaystyle \int \frac{\dd x}{x^2 - 6x + 25}$ and
$\displaystyle \int \frac{2x - 1}{x^2 - 6x + 25}\,\dd x.$ For the second, split
the numerator into a multiple of the denominator's derivative plus a constant.
\end{exercise}

\begin{exercise}[Trig substitution]
Integrate $\displaystyle \int \frac{x^2}{\sqrt{4 - x^2}}\,\dd x$ using the
sine triangle. Draw the triangle, convert every quantity, integrate the
resulting trig expression, and convert your answer back to $x$ explicitly.
\end{exercise}

\begin{exercise}[Trig substitution, the other two triangles]
Integrate $\displaystyle \int \frac{\dd x}{x^2\sqrt{x^2 - 9}}$ using the secant
triangle, and $\displaystyle \int \sqrt{x^2 + 1}\,\dd x$ using the tangent
triangle. The second will need $\int \sec^3\theta\,\dd\theta$; derive it with
one integration by parts.
\end{exercise}

\begin{exercise}[Conceptual / proof]
Prove that the partial-fraction building blocks $\dfrac{Q(x)}{x - r_k}$, where
$Q(x) = \prod_{i=1}^n (x - r_i)$ has $n$ distinct roots, are linearly
independent as polynomials. (Mimic the cover-up argument of
Section~\ref{ch24:why-pf-exists}.) Explain in a sentence why this independence
is exactly what guarantees the partial-fraction decomposition exists for every
proper numerator.
\end{exercise}

\begin{exercise}[Synthesis]
The single integral $\displaystyle \int \frac{\dd x}{\sqrt{x^2 + 2x + 2}}$ uses
two techniques from this chapter in sequence. Name them in order, carry the
computation through to a function of $x$, and write one sentence explaining how
this problem embodies the chapter's slogan that every method reshapes the
integrand into a known form.
\end{exercise}
